\section{Critical points and an equivariant pairing}
\label{sec:critical-pairing}

For targets with the compact superlevels obtained in the preceding section,
we can compare the critical points of $V^F_\xi$.  We will assign every saddle to a distinct
maximum whose value is at least as large.  The assignment must depend
measurably on $F$ and respect affine changes of the source, because the next
section will average this comparison using the weighted law.

\subsection{The critical map and good targets}

Throughout this section, $P$ is the law of Lemma~\ref{lem:weighted-law},
$\beta>1$, and $k=(\beta+3)/4>1$.  Recall that
$Q_F=1/F'$ and $q_F(x+iy)=iyQ'_F(x+iy)/Q_F(x+iy)=a_F+ib_F$.
Define the smooth real map $G_F:\HH\to\CC$ by
\begin{equation}
 G_F(z)=F(z)-\frac{iy}{k}F'(z),\qquad z=x+iy.
 \label{eq:critical-map}
\end{equation}
The normalization gives $G_F(i)=-i/k$ for every $F\in\KK$.
We define its normalized signed Jacobian by
\begin{equation}
 J_F(z)=\frac{\det DG_F(z)}{|F'(z)|^2}.
 \label{eq:normalized-jacobian}
\end{equation}
Here $DG_F$ is the derivative as a map between two real planes.  The sign
of $J_F$ will distinguish saddles from maxima.

\begin{lemma}[Jacobian and its expectation]
\label{lem:jacobian-moments}
For every $F\in\KK$ and $z\in\HH$,
\begin{equation}
 k^2J_F(z)=k(k-1)+(2k-1)a_F(z)+|q_F(z)|^2.
 \label{eq:jacobian-formula}
\end{equation}
The function $F\mapsto J_F(i)$ is bounded and continuous on $\KK$, and
\begin{equation}
 k^2\EE J_F(i)=k(k-1)>0.
 \label{eq:positive-jacobian-mean}
\end{equation}
\end{lemma}

\begin{proof}
Since $F''/F'=iq_F/y$, the two real derivative columns are
\[
 (G_F)_x=\frac{F'}{k}(k+q_F),\qquad
 (G_F)_y=\frac{iF'}{k}(k-1+q_F).
\]
For complex vectors $u,v$, their real determinant is
$\operatorname{Im}(\overline u v)$.  Applying this identity to the displayed
columns gives \eqref{eq:jacobian-formula}.  Local derivative continuity and
compactness of $\KK$ give the asserted continuity and boundedness at $i$.
Lemma~\ref{lem:law-moments} gives
\[
 \EE a_F(i)=-\frac\beta2,\qquad
 \EE|q_F(i)|^2=\frac{\beta(\beta+1)}4.
\]
Because $2k-1=(\beta+1)/2$, the last two terms in the expectation of
\eqref{eq:jacobian-formula} cancel.  This proves
\eqref{eq:positive-jacobian-mean}.
\end{proof}

\begin{lemma}[Critical points and their signs]
\label{lem:critical-signs}
Fix $F\in\KK$ and $\xi\in\CC$.  Away from the unique possible pole of
$V^F_\xi$, put $\ell=\log V^F_\xi$.  Then
\begin{equation}
 \Delta\ell=-\frac{k}{y^2},
 \label{eq:negative-laplacian}
\end{equation}
and its critical points are exactly the solutions of $G_F(z)=\xi$.
At such a point,
\begin{align}
 V^F_\xi(z)&=k y^{k-1}|Q_F(z)|,
 \label{eq:critical-height}\\
 \operatorname{Hess}\ell(z)
 &=\frac{k}{y^2}
 \begin{pmatrix}
  -(k+a_F(z))&b_F(z)\\
  b_F(z)&k-1+a_F(z)
 \end{pmatrix},
 \label{eq:critical-hessian}\\
 \det\operatorname{Hess}\ell(z)&=-\frac{k^4}{y^4}J_F(z).
 \label{eq:critical-hessian-det}
\end{align}
Consequently a solution with $J_F(z)>0$ is a nondegenerate saddle,
and one with $J_F(z)<0$ is a nondegenerate maximum.
\end{lemma}

\begin{proof}
Write $r=F'/(F-\xi)$.  Since $\log|F-\xi|$ is harmonic off the pole,
\[
 \ell_x=-\operatorname{Re}r,\qquad
 \ell_y=\frac{k}{y}+\operatorname{Im}r,
\]
and \eqref{eq:negative-laplacian} follows.  Thus the critical equation is
$r=-ik/y$, or $F-\xi=iyF'/k$, which is precisely $G_F(z)=\xi$ and gives
\eqref{eq:critical-height}.  A pole cannot solve this equation because
$y>0$ and $F'\ne0$.

At a critical point the holomorphic derivative of $r$ is
\[
 r'=\frac{F''}{F'}r-r^2
     =\frac{k(k+q_F)}{y^2}.
\]
Differentiating the two real gradient entries now gives
\eqref{eq:critical-hessian}, whose determinant is
\eqref{eq:critical-hessian-det} by \eqref{eq:jacobian-formula}.
Its trace is $-k/y^2<0$.  A negative determinant therefore gives a saddle;
a positive determinant gives two negative eigenvalues and hence a strict
maximum.  At a critical point,
$\operatorname{Hess}V^F_\xi=V^F_\xi\operatorname{Hess}\ell$, so these
classifications apply to $V^F_\xi$ as well.
\end{proof}

\begin{definition}
\label{def:good-pair}
A pair $(F,\xi)\in\KK\times\CC$ is \emph{good} if every closed positive
superlevel of $V^F_\xi$ is compact in $\HH$, and $\xi$ is a regular value
of $G_F$.  Thus every solution $G_F(z)=\xi$ must satisfy $J_F(z)\ne0$;
a target with no such solution is regular.  Write $\mathcal G$ for the
set of good pairs.
\end{definition}

\begin{lemma}[Good pairs and affine changes of source]
\label{lem:good-pairs}
The set $\mathcal G$ is Borel.  For $P$-almost every $F$, its section
$\{\xi:(F,\xi)\in\mathcal G\}$ has full planar area measure.

For every $F\in\KK$, $\xi\in\CC$, $t=x_0+iy_0\in\HH$, and $w\in\HH$, set
$\widetilde\xi=(\xi-F(t))/(y_0F'(t))$.  Then
\begin{align}
 G_{T_tF}(w)&=\frac{G_F(t\cdot w)-F(t)}{y_0F'(t)},
 \label{eq:critical-map-covariance}\\
 q_{T_tF}(w)&=q_F(t\cdot w),\qquad J_{T_tF}(w)=J_F(t\cdot w),
 \label{eq:jacobian-covariance}\\
 V^{T_tF}_{\widetilde\xi}(w)
 &=y_0^{1-k}|F'(t)|V^F_\xi(t\cdot w).
 \label{eq:critical-height-covariance}
\end{align}
In particular, $(F,\xi)\in\mathcal G$ if and only if
$(T_tF,\widetilde\xi)\in\mathcal G$.
\end{lemma}

\begin{proof}
The compact-superlevel condition is Borel by
Lemma~\ref{lem:superlevel-borel}.  For completeness, regularity can be tested
on the compact source boxes
$C_m=\{|x|\le m,\ 1/m\le y\le m\}$: failure means that for some $m$
there is $w\in C_m$ with
\[
 G_F(w)=\xi,\qquad J_F(w)=0.
\]
The set defined by these equalities is closed in
$\KK\times\CC\times C_m$; its projection onto $\KK\times\CC$ is closed
because $C_m$ is compact.  A countable union gives a Borel exceptional set.

For each fixed $F$, the critical values of $G_F$ have area zero. This is
the plane case of Sard's theorem \cite[Theorem~4.1]{Sard1942}; we give
the short covering argument for the smooth maps at hand.  Fix a compact source box and a slightly
larger compact box in $\HH$.  For small $h>0$, cover the first box by
$O(h^{-2})$ squares of side $h$.  If a square meets the critical set, choose
a critical point there.  Taylor expansion at that point, using the bounded
second derivative on the larger box, puts the image of the square in an
$O(h^2)$-neighborhood of a line segment of length $O(h)$: the linear part
has rank at most one.  This neighborhood has area $O(h^3)$.
The image of all critical points in the first box consequently has outer
area $O(h)$, and therefore area zero.  Exhausting $\HH$ by boxes proves the
claim.  Combine it with Proposition~\ref{prop:proper-superlevels} to obtain
the asserted full-area sections for almost every $F$.

The derivative identity
$(T_tF)'(w)=F'(t\cdot w)/F'(t)$ gives
\eqref{eq:critical-map-covariance} by substitution and gives the $q$
identity in \eqref{eq:jacobian-covariance}.  Formula
\eqref{eq:jacobian-formula} then gives the $J$ identity.
The definition of $V$ gives \eqref{eq:critical-height-covariance}.
Its multiplying factor is positive and independent of $w$.
The source affine map is a homeomorphism of $\HH$, while the target change
is a nonsingular similarity.  They preserve the two conditions defining
$\mathcal G$ in both directions.
\end{proof}

\subsection{The superlevel count}

We next prove the counting statement needed for the pairing.  The pole is
kept separate from the finite maxima: its contribution to the count is an
index of the gradient, rather than a critical point of the smooth function
$\ell$.

\begin{lemma}[Counting critical points above a level]
\label{lem:critical-count}
Fix a good pair $(F,\xi)$.  There are only finitely many critical points
of $V^F_\xi$ with value at least $h$, for every $h>0$.
For every such $h$,
\begin{equation}
 \#\{\text{maxima with }V^F_\xi\ge h\}
 \ \ge\
 \#\{\text{saddles with }V^F_\xi\ge h\}.
 \label{eq:cumulative-critical-count}
\end{equation}
The possible pole is excluded from the maximum count.
\end{lemma}

\begin{proof}
Put $\ell=\log V^F_\xi$ away from the pole.
The solutions of $G_F(z)=\xi$ are isolated by regularity, and form a closed
subset of $\HH$; isolation follows from the real inverse-function theorem
\cite[Theorem~C.34]{Lee2013}.  Those with $V^F_\xi\ge h$ lie in a compact set and are
therefore finite.  This includes the possibility of a pole in that compact
set without allowing accumulation there: if $F(z_0)=\xi$, then
$G_F(z_0)-\xi=-i\operatorname{Im}(z_0)F'(z_0)/k\ne0$.

Choose $\lambda>0$ which is not a critical value of $V^F_\xi$ and put
\[
 U_\lambda=\{z\in\HH:V^F_\xi(z)>\lambda\},
\]
including the pole if it exists.  This is open, with compact closure in
$\HH$.  Its boundary is exactly the level $\{V^F_\xi=\lambda\}$:
the pole is an interior point of $U_\lambda$, and at every level point the
nonzero gradient gives values on both sides of $\lambda$.
The level is a compact smooth one-manifold without boundary.  Its
components are open in that manifold, and thus there are finitely many.
Each component is a smooth regular simple closed curve.  One way to see
the last assertion is to flow a unit tangent field, obtained by rotating
$\nabla\ell/|\nabla\ell|$ along the level.  The flow is complete by
compactness.  Its orbits are open in each connected component, so a
component is a single orbit.  If this orbit were not periodic, its
parametrization would be an injective surjective local homeomorphism from
$\mathbb R$ to that compact component, which is impossible.  The return
times form a closed subgroup with a gap around zero by local injectivity;
its least positive element gives a regular simple closed parametrization.

Consider the bounded Jordan interior $D$ of one level curve $\gamma$.
The Jordan curve theorem is used here in its plane separation form
\cite[Theorem~63.4]{Munkres2014}.  The closure of $D$ lies in $\HH$:
points below the positive minimum height of $\gamma$ connect to infinity
without crossing it and hence lie in the unbounded complementary
component.  In fact
\begin{equation}
 V^F_\xi>\lambda\quad\hbox{throughout }D.
 \label{eq:jordan-interior-superlevel}
\end{equation}
If a point inside had smaller value, $\ell$ would attain a minimum below
its boundary value $\log\lambda$ at an interior point.  A neighborhood of
any pole can be excluded when taking this minimum, since $\ell$ tends to
$+\infty$ there.  Such a minimum contradicts
$\Delta\ell=-k/y^2<0$.  Once smaller values are excluded, equality at an
interior point is likewise a local minimum and is impossible.

The interiors of the level curves are disjoint and are exactly the
components of $U_\lambda$.  Indeed, nested curves are excluded by
\eqref{eq:jordan-interior-superlevel}.  Conversely, a nonempty component
of $U_\lambda$ has a boundary point because its closure is compact.
This point is on the level: a boundary point cannot belong to a component
of the open set $U_\lambda$.  Near it, regularity and
\eqref{eq:jordan-interior-superlevel} identify the superlevel side with
the interior side of its level curve.  The component therefore meets and
contains that connected Jordan interior, and cannot leave it without
crossing the curve.  This proves the assertion, including that no
component has a hole.

The superlevel components are therefore bounded Jordan domains with smooth boundaries.
It remains to turn this topological information into a count of their
critical points, keeping the possible pole visible in the index sum.
On the positive orientation of each boundary curve, $\nabla\ell$ is an
inward normal and has winding number $+1$.  We recall the elementary
tangent-turning argument of Hopf~\cite[\S2, pp.~53--55]{Hopf1935} to fix this sign.
Parametrize a regular simple closed curve periodically by
$\gamma:[0,1]\to\mathbb R^2$, positively oriented and starting at a
leftmost point.  Its initial tangent points down.  The unit secant
\[
 \frac{\gamma(t)-\gamma(s)}{|\gamma(t)-\gamma(s)|},
 \qquad 0\le s<t\le1,\quad(s,t)\ne(0,1),
\]
extends continuously to the closed triangle: on the diagonal its value
is the unit tangent, and at $(0,1)$ it is the negative initial tangent.
The latter extension follows by the endpoint expansion
\[
 \gamma(t)-\gamma(s)
 =-\bigl((1-t)+s\bigr)\gamma'(0)
   +o\bigl((1-t)+s\bigr).
\]
The extended secant has a continuous angle lift $A(s,t)$ on the triangle.
For example, uniform continuity permits one common number of subdivisions
of every segment from $(0,0)$ to $(s,t)$ so that successive unit-vector
ratios have a unique small continuous argument; adding those arguments
constructs the lift.  Along $(0,t)$, secants lie in the closed right
semicircle, from the downward tangent to its upward negative.  Thus
$A(0,1)-A(0,0)=\pi$.  The secant along $(t,1)$ is the negative of that
along $(0,t)$, including their extensions, so continuity gives
$A(t,1)-A(0,t)=\pi$ throughout.  Therefore
$A(1,1)-A(0,0)=2\pi$.  This is the turn of the tangent along the diagonal;
rotation to the inward normal preserves its winding number.

Inside one component of $U_\lambda$, remove small disjoint disks about its
finitely many critical points and its pole, if present.  On the remaining
compact smooth domain, $X=\nabla\ell$ is nonzero and its angle form
\[
 \frac{X_1\dd X_2-X_2\dd X_1}{X_1^2+X_2^2}
\]
is closed.  Stokes' theorem~\cite[Theorem~16.11]{Lee2013} shows that the
winding on the outer boundary equals the sum of the windings on the
positively oriented small circles.  At a nondegenerate critical point,
the gradient on a sufficiently small circle is homotopic through
nonzero vectors to its invertible linear part.  Its index is consequently
the sign of $\det\operatorname{Hess}\ell$, namely $+1$ at a maximum and
$-1$ at a saddle.

If $z_0$ is the pole, the simple zero of $F-\xi$ gives
\[
 \ell(z)=-\log|z-z_0|+\text{a smooth function},\qquad
 \nabla\ell(z)=-\frac{z-z_0}{|z-z_0|^2}+O(1),
\]
where vectors are identified with complex numbers in the second formula.
The negative radial leading field has winding number $+1$.  Thus the pole
also contributes $+1$, although it is not a finite critical point.

Let $C_\lambda$ be the number of components of $U_\lambda$, and let
$B\in\{0,1\}$ be the number of poles.  Summing the index identity gives
\begin{equation}
 C_\lambda
 =\#\{\text{maxima with }V^F_\xi>\lambda\}
  -\#\{\text{saddles with }V^F_\xi>\lambda\}+B.
 \label{eq:level-index-count}
\end{equation}
If $B=1$, its pole belongs to $U_\lambda$, so $C_\lambda\ge B$.
The same inequality is automatic when $B=0$.  Hence maxima are at least
as numerous as saddles above every regular positive level.  An empty
superlevel gives the same identity with all counts zero.

Finally, for a given $h>0$, finiteness of the critical points of value at
least $h/2$ permits a regular $\lambda\in(h/2,h)$ with no critical value
in $[\lambda,h)$.  The counts above $\lambda$ are exactly those at or
above $h$, including all ties.  This proves
\eqref{eq:cumulative-critical-count}.
\end{proof}

\subsection{Equal-rank pairing}

The cumulative count supplies the partner of each saddle.  We now give
one choice which is intrinsic to the source coordinates and prove that
it has the measurability needed for integration.

\begin{proposition}[Measurable equivariant pairing]
\label{prop:critical-pairing}
Let
\[
 \mathcal D=\{(F,z)\in\KK\times\HH:
                  J_F(z)>0,\ (F,G_F(z))\in\mathcal G\}.
\]
There is a Borel map $R:\mathcal D\to\HH$, written $R_F(z)=R(F,z)$,
such that, for every $(F,z)\in\mathcal D$, with $z'=R_F(z)$ and
$\xi=G_F(z)$,
\begin{equation}
 G_F(z')=\xi,\qquad J_F(z')<0,\qquad
 V^F_\xi(z')\ge V^F_\xi(z).
 \label{eq:pairing-guarantees}
\end{equation}
For each fixed $F$, the map $R_F$ is injective on its domain.
It is affine equivariant: if $t\in\HH$ and $(F,t\cdot w)\in\mathcal D$, then
\begin{equation}
 (T_tF,w)\in\mathcal D,\qquad
 R_{T_tF}(w)=t^{\#}\cdot R_F(t\cdot w).
 \label{eq:pairing-covariance}
\end{equation}
\end{proposition}

\begin{proof}
Fix a good pair $(F,\xi)$.  Order its saddles by decreasing value of
$V^F_\xi$, breaking ties by increasing lexicographic order of $(x,y)$.
Order its maxima separately by the same rule.  Every point has finitely
many predecessors, since all predecessors have value at least its own
positive value.  Give it the zero-based rank equal to the number of its
predecessors.  Distinct points have distinct ranks: if one precedes
another, its predecessors and the point itself precede the latter.
Moreover, if rank $j$ occurs, the finite ordered set of its predecessors
has ranks $0,\ldots,j-1$.  Thus there are no missing ranks below an
attained one, whether the entire fiber is finite or infinite.

Suppose a saddle has rank $j$ and value $v$.  There are at least $j+1$
saddles of value at least $v$.  Lemma~\ref{lem:critical-count} gives at
least $j+1$ maxima of value at least $v$.  These maxima form an initial
segment of the maximum order, so the maximum of rank $j$ exists and has
value at least $v$.  Assign the saddle to this maximum.  This is
injective within the fiber, and gives \eqref{eq:pairing-guarantees} by
Lemma~\ref{lem:critical-signs}.  Distinct fibers cannot share a point,
so the assignment is injective for fixed $F$ across its entire domain.

Figure~\ref{fig:rank-pairing} illustrates the two ordered lists.
\begin{figure}[htbp]
\centering
\begin{tikzpicture}[x=1.05cm,y=.72cm,>=Latex,
  every node/.style={font=\small}]
 \draw[->] (-1.1,0)--(-1.1,4.3) node[above] {$V^F_\xi$};
 \node at (0,4.6) {saddles};
 \node at (4,4.6) {maxima};
 \foreach \sy/\my/\r in {3.5/4/0,2.7/3.1/1,.6/1.4/2} {
  \draw[->,thick] (.09,\sy)--(3.9,\my);
  \fill (0,\sy) circle (2pt) node[left] {$s_{\r}$};
  \fill (4,\my) circle (2pt) node[right] {$m_{\r}$};
 }
\end{tikzpicture}
\caption{The first three ranks in a schematic good target fiber. Vertical
position records $V^F_\xi$; the two columns are ordered separately by decreasing
value, with lexicographic tie-breaking. Equal ranks are paired, and each arrow
ends at least as high as it starts.
The pole is absent from both lists, and additional maxima may remain unassigned.
Arrows describe the rank assignment, not curves in the source plane.}
\label{fig:rank-pairing}
\end{figure}


We give the Borel construction explicitly. Countable Borel enumeration
is the subject of the Lusin--Novikov theorem; see
\cite[4F.6 and 4F.17]{Moschovakis2009}. Here nonsingularity supplies
continuous local branches, and the rank definition will also supply
affine equivariance. Consider the nonsingular
solution graph
\[
 \Gamma=\{(F,\xi,w)\in\KK\times\CC\times\HH:
                    G_F(w)=\xi,\ J_F(w)\ne0\}.
\]
It can be covered by countably many continuous branches over open
parameter sets in $\KK\times\CC$.  To prove this without differentiating
in the parameter $F$, fix $(F_0,\xi_0,w_0)\in\Gamma$ and the invertible
real matrix $D=DG_{F_0}(w_0)$.  Joint continuity of derivatives on source
compacts gives a closed ball $\overline B(w_0,r)\subset\HH$ and an open
parameter neighborhood $U$ such that
\[
 \|I-D^{-1}DG_F(w)\|\le\frac12
 \quad(w\in\overline B(w_0,r)),\qquad
 |D^{-1}(G_F(w_0)-\xi)|<\frac r2
\]
for every $(F,\xi)\in U$.  The map
\[
 w\longmapsto w-D^{-1}(G_F(w)-\xi)
\]
is a contraction of the closed ball into its interior.  Its unique fixed
point exists by the contraction principle~\cite[Lemma~C.35]{Lee2013}.  This
point solves $G_F(w)=\xi$, is nonsingular, and depends continuously on
$(F,\xi)$.  Continuity follows, for instance, by taking any convergent
parameter sequence and using compactness of the ball and uniqueness to
identify every subsequential limit of its fixed points.

Consequently $\Gamma\cap(U\times B(w_0,r))$ is the graph of one
continuous branch over $U$, and is an open neighborhood in $\Gamma$.
The space $\Gamma$ is second countable, being a subspace of
$\KK\times\CC\times\HH$.  A countable collection of these branch
neighborhoods therefore covers it.  Enumerate the resulting continuous
partial branches as $b_j:U_j\to\HH$, $j\ge1$.
Activate branch $j$ at a parameter if the parameter is in $U_j$ and its
value differs from every earlier branch defined there.  These activation
sets are Borel: there are finitely many earlier indices for a given $j$,
and equality of two continuous branches is closed relative to their
open overlap.  Each nonsingular root occurs in exactly one active branch.

On a branch, its critical value can be evaluated without division by
$F(w)-\xi$, using the positive continuous function
\[
 H(F,w)=k(\operatorname{Im}w)^{k-1}|Q_F(w)|.
\]
At every root this equals $V^F_\xi(w)$ by
\eqref{eq:critical-height}.  The signs of $J_F$, comparisons of $H$,
equalities of values, and lexicographic comparisons of coordinates are
Borel.  The rank of an active branch among branches of its sign is thus
the countable sum of the Borel indicators of its active predecessors.
This is an extended-integer-valued Borel function, and is finite on good
fibers.  Requiring a positive-sign and a negative-sign active branch to
have the same finite rank gives a countable Borel test for the assigned
pair.

The set $\mathcal D$ is Borel by Lemma~\ref{lem:good-pairs} and continuity
of $(F,z)\mapsto(F,G_F(z))$.  At this target, require additionally that
the positive-sign branch equal $z$.  Exactly one active negative-sign
branch passes the equal-rank test, by the existence and uniqueness proved
above.  Selecting its value defines a Borel map on $\mathcal D$: the
preimage of an open source set is the countable union of the Borel tests
just described with the selected branch constrained to that set.

Finally, under the affine source change $w\mapsto t\cdot w$, the critical
fibers and signs correspond by
\eqref{eq:critical-map-covariance}--\eqref{eq:jacobian-covariance}.
All critical heights in a fiber are multiplied by the same positive
constant in \eqref{eq:critical-height-covariance}.  The positive real
source scale preserves lexicographic order, including ties of the first
coordinate.  Hence the two orders, their ranks, and the equal-rank
assignment correspond exactly, giving \eqref{eq:pairing-covariance}.
The enumeration was used only to prove measurability and does not enter
the geometric definition of the assignment.
\end{proof}

The ordering used in \eqref{eq:pairing-guarantees} also handles coincident critical values.
At each critical point, \eqref{eq:critical-height} and
$4(k-1)=\beta-1$ give
\begin{equation}
 y^{\beta-1}|Q_F(z)|^4
   =\left(\frac{V^F_\xi(z)}k\right)^4.
 \label{eq:critical-transport-weight}
\end{equation}
Thus the pairing increases the weight on the left of
\eqref{eq:critical-transport-weight}.  The next section uses this
pointwise comparison to contradict the positive expectation in
\eqref{eq:positive-jacobian-mean}.
